\documentclass[11pt]{article}
\usepackage[T1]{fontenc}
\usepackage{lmodern,amsmath,amssymb,amsthm,mathtools,booktabs,array}
\usepackage[margin=1in]{geometry}
\usepackage[colorlinks=true,linkcolor=blue,citecolor=blue,urlcolor=blue]{hyperref}
\usepackage{microtype}
\newtheorem{theorem}{Theorem}[section]
\newtheorem{lemma}[theorem]{Lemma}
\newtheorem{proposition}[theorem]{Proposition}
\newtheorem{remark}[theorem]{Remark}
\newcommand{\Z}{\mathbb Z}
\newcommand{\R}{\mathbb R}
\newcommand{\conv}{\operatorname{conv}}
\newcommand{\area}{\operatorname{area}}
\newcommand{\intr}{\operatorname{int}}
\newcommand{\width}{\operatorname{width}}
\newcommand{\Ehr}{\mathcal E}
\title{Counting Ehrhart functions of half-integral polygons}
\author{Kaiwen Jin\\\small\href{mailto:yui.ui.siki@gmail.com}{yui.ui.siki@gmail.com}}
\date{}
\begin{document}
\maketitle
\begin{abstract}
Fix an integer \(i\ge2\). We count the distinct Ehrhart functions of convex
half-integral polygons with exactly \(i\) interior lattice points and at least
two boundary lattice points. The number is
\((9i^3+72i^2+175i+106)/2\), as proposed by Bohnert and Springer.
We describe the complete set of functions by their area and their boundary
counts before and after doubling. Convex section estimates isolate two
exceptional families, and explicit polygons realize every remaining
signature. The proof uses the published complete computational classification
for \(2\le i\le16\), five published minimum-width bounds, and a finite list of
exact rational inequality certificates. Integral polygons and period collapse
are included, and no restriction is imposed on the arrangement of the
interior lattice points.
\end{abstract}

\section{Introduction}
For a rational polygon \(P\), its Ehrhart function records the lattice points
in every integer dilate:
\[
 L_P(t)=\#(tP\cap\Z^2),\qquad t\in\Z_{\ge0}.
\]
For lattice polygons, Pick's formula determines this polynomial from the
interior and boundary counts. Scott's inequality bounds the possible boundary
counts when the interior count is fixed \cite{Scott}. Thus counting the
functions is much simpler than classifying the polygons themselves.

For half-integral polygons, meaning polygons whose vertices lie in
\(\frac12\Z^2\), there are two polynomial constituents. The boundary count of
the doubled polygon supplies an additional parameter. Bohnert's coefficient
formula makes these parameters explicit \cite[Proposition 2.1]{Collapse}.
Bohnert and Springer established sharp boundary and area bounds
\cite[Theorems 1.1 and 1.4]{Inequalities}; these bounds constrain the
coefficients but do not determine which intermediate values of the doubled
boundary count occur. Their published classification enumerates all functions
through \(i=16\) and proposes a cubic formula for the number with at least
two boundary lattice points \cite[Conjecture 6.8]{Classification}.

We prove that formula.
\begin{theorem}\label{thm:count}
For every integer \(i\ge2\), the number of distinct functions \(L_P\) among
all two-dimensional convex polygons \(P\) with vertices in
\(\frac12\Z^2\), exactly \(i\) interior \(\Z^2\) points and at least two
boundary \(\Z^2\) points is
\[
 \Ehr(i)=\frac{9i^3+72i^2+175i+106}{2}.
\]
\end{theorem}
Here half-integral includes integral polygons. Distinctness means equality
of the entire function on \(\Z_{\ge0}\), so a function with period collapse
is counted only once. The class is exactly that in the published conjecture;
in particular, the interior points need not be collinear.

There are two geometric obstructions to a direct count from the marginal
area bounds. First, a polygon can have an integer hull which is a segment.
A thin family then has a smaller doubling interior count and must be treated
separately. Second, the upper bound coupling area and doubled boundary count
has a width-two equality family. We identify both families. Outside them,
section estimates give a finite set of inequalities. Uniform constructions
then show that those inequalities describe all signatures, rather than only
necessary restrictions.

The proof retains explicit computational inputs. Small interior counts use
the complete published classification, and a near-area estimate uses five
minimum-width maxima from the same maximal-polygon catalogue. The wide
section inequality is proved by rational linear combinations for finitely
many width indices and by an analytic bound for the remaining indices.
The all-parameter constructions and the count itself are proved
symbolically. These roles are distinguished in Section~\ref{sec:computation}.

We first give the signature classification, then prove its necessary bounds
and realize every signature. The final count is in Section~\ref{sec:count}.
Throughout, area is Euclidean area, lattice points without a qualifier belong
to \(\Z^2\), and all polygons are convex and two-dimensional.

\section{Signatures and the classification}\label{sec:signatures}
Write
\[
 i=i(P),\quad b=b(P),\quad Q=2P,\quad
 A=8\area(P)=2\area(Q),\quad B=b(Q).
\]
Thus \(A,B\) are integers. Bohnert's coefficient formula
\cite[Proposition 2.1]{Collapse} gives
\begin{equation}\label{eq:coeff}
 L_P(t)=
 \begin{cases}
 \dfrac A8t^2+\dfrac b2t+i+\dfrac b2-\dfrac A8,&t\ \text{odd},\\[2mm]
 \dfrac A8t^2+\dfrac B4t+1,&t\ \text{even}.
 \end{cases}
\end{equation}
The even formula includes \(t=0\). At fixed \(i\), different triples
\((A,b,B)\) give different functions: the even constituent determines \(A,B\)
and the odd linear coefficient determines \(b\).
Period collapse occurs precisely when
\begin{equation}\label{eq:collapse}
 B=2b,\qquad A=8i+4b-8.
\end{equation}

Define three sets of integer triples. For the \emph{body}, let
\(2\le b\le3i+6\), put
\[
 q(b)=
 \begin{cases}8,&b\le3i+4,\\7,&b=3i+5,\\6,&b=3i+6,\end{cases}
\]
and require
\begin{align}
 6i+4b-6&\le A\le12i+2b+q(b),\label{eq:bodyA}\\
 A&\equiv B\pmod2,\label{eq:parity}\\
 2b+(A\bmod2)+2\delta&\le B
 \le\min\{A-6i-2b+6,\ 3i+b+6,\ 12i+4b+10-A\}.\label{eq:bodyB}
\end{align}
Here \(\delta=1\) if \(A=12i+2b+8\) and \(b\ne2i+4\), or if
\(b=2,A=12i+10\); otherwise \(\delta=0\).
The \emph{top} is
\begin{equation}\label{eq:topset}
 (A,b,B)=(12i+2b+8-d,\ b,\ 2b+4+d),\quad
 d\in\{0,1,2\},\quad b\ge2,\quad b+2d\le2i+2.
\end{equation}
The \emph{line} is
\begin{equation}\label{eq:lineset}
 (A,b,B)=(A,2,A-4i),\qquad 4i+4\le A\le8i+8.
\end{equation}
Let \(C_i\) be their union.
\begin{theorem}\label{thm:signatures}
For every \(i\ge2\), the signatures \((A,b,B)\) of the polygons in
Theorem~\ref{thm:count} are exactly \(C_i\).
\end{theorem}
We prove necessity for \(i\ge17\) in the next two sections. For
\(2\le i\le16\) we use the complete published classification as specified in
Section~\ref{sec:computation}. Section~\ref{sec:witnesses} proves attainment
for all \(i\ge2\).

We use the following published marginal bounds. For a nonintegral
half-integral polygon with \(i\ge2\),
\begin{equation}\label{eq:marginal}
 b\le3i+6,\qquad A\le12i+2b+q(b)
\end{equation}
by \cite[Theorems 1.1 and 1.4]{Inequalities}. At \(b\ge3\), its
Proposition 3.3 also gives
\begin{equation}\label{eq:maxlower}
 B\ge2b+2\quad\text{if }A=12i+2b+8,\ b\ne2i+4.
\end{equation}
Only this necessary assertion is used from that proposition.
For integral polygons, Pick and Scott instead give
\(A=8i+4b-8,\ B=2b,\ b\le2i+6\); these satisfy the marginal bounds
and cannot reach the exceptional maximum area in \eqref{eq:maxlower}.
In general \(B\ge2b\): between successive original boundary points is a
nonoriginal half-grid boundary point. Pick on \(Q\) gives
\begin{equation}\label{eq:pickQ}
 A=2i(Q)+B-2,\qquad L_P(2)=\frac{A+B+2}{2},
\end{equation}
and hence \eqref{eq:parity}.

\section{Doubling and convex sections}\label{sec:bounds}
\subsection{The lower doubling bound}
\begin{lemma}\label{lem:lower}
Either \(i(Q)\ge3i+b-2\), or the signature belongs to \eqref{eq:lineset}.
\end{lemma}
\begin{proof}
Let \(S=P\cap\Z^2\), \(K=\conv(S)\). If \(K\) is two-dimensional,
triangulate it with every point of \(S\) as a vertex. All edges are primitive.
Writing \(h=i(K)\), the number of edges is
\[
 E=2i+2b+h-3.
\]
This follows from the fan triangulation formula
\(E=3\#S-b(K)-3\), since \(b(K)=i+b-h\); successive insertions preserve it.

Discard boundary edges whose two endpoints are boundary points of \(P\).
All other edge interiors lie in \(\intr(P)\). Their primitive midpoints are
nonintegral and distinct, so they yield distinct interior points of \(Q\),
in addition to the \(i\) doubled original interior points.
If \(s\) edges were discarded, \(i(Q)\ge i+E-s\).
For \(h\ge1\), the cyclic boundary has at most \(b\) discarded edges.
For \(h=0\), it contains both original interior and boundary types, so a
boundary-type run of length \(r\) contributes only \(r-1\) discarded edges;
thus \(s\le b-1\). Both give the required bound.

If \(K\) is a segment, an integer affine unimodular transformation puts its
line at \(y=0\). Since \(P\) has interior points on that line, it extends
on both sides. All its original boundary points are the two intersection
endpoints. Thus \(b=2\) and we may take them as \((0,0),(R,0)\), \(R=i+1\).
A vertex of height \(H\ge3/2\) would give at \(y=1\) a triangle section of
length \(R(1-1/H)\ge1\), containing an off-line original point.
The reflected argument excludes heights at most \(-3/2\).

If a vertex has height \(1\), its horizontal coordinate \(x\) is a
nonintegral half-integer. Its doubled triangle contains \(R\) interior
fine points on \(y=1\), while the middle row contains \(2R-1\), so
\(i(Q)\ge3R-1=3i+2\). The same holds at height \(-1\).
Otherwise all vertices have heights in \(\{-1/2,0,1/2\}\).
Only the middle fine row is interior and \(i(Q)=2i+1\).
The two triangles over the middle base give \(A\ge4R\).
If \(\ell_-,\ell_+\) are the original end-section lengths, concavity gives
\(\ell_-+\ell_+\le2R\), and
\(\area(P)=(\ell_-+2R+\ell_+)/4\le R\).
Pick now gives \(B=A-4i\), proving the line alternative.
\end{proof}
Outside that line, \eqref{eq:pickQ} gives
\(B\le A-6i-2b+6\); with \(B\ge2b\), this also gives the lower
area bound in \eqref{eq:bodyA}.

\subsection{Minimizing sections and the boundary cap}
Minimum lattice width is the minimum of
\(\max_P\langle u,x\rangle-\min_P\langle u,x\rangle\) over nonzero
primitive integer covectors \(u\). It is attained: an inscribed ball of
radius \(r>0\) gives width at least \(2r|u|\), so only finitely many
integer directions can improve a fixed direction.
An integer unimodular map puts a minimizer on a coordinate axis.

Bohnert's section estimate \cite[Lemma 2.12]{AreaBounds} implies, in that
coordinate, that a section of a convex body of width \(w\) has length at least
\(\min\{w/2,\text{distance from the nearer projection end}\}\).
This follows from either of his two bounds relative to the longest section.
It is preserved by doubling.

Index the fine columns of \(Q\) by \(0,\ldots,W\). An E end has even
original fine horizontal coordinate; an O end has odd coordinate.
This is only an indexing convention. Original points still have even fine
vertical coordinate, and lie on the parity of columns prescribed by E/O.
Write \(u,v\) for the upper/lower ordinates, \(\ell=u-v\).
For one concave chain \(f=u\) or \(f=-v\), let \(\beta_j\) indicate an
integer ordinate and \(c_j\) an even integer ordinate on a coarse column.
On an interior coarse column,
\[
 i_j=\left\lceil u_j/2\right\rceil+
       \left\lceil-v_j/2\right\rceil-1,\qquad \ell_j\le2i_j+2.
\]

\begin{lemma}\label{lem:boundarycap}
For \(i\ge13\), a signature outside the line family satisfies \(B-b\le3i+6\).
\end{lemma}
\begin{proof}
Write \(N=B-b\). The two narrow nonthin widths are \(W=3,4\).
For \(W=3\), reflect to pattern EO, with coarse interior column 2.
Put \(M=\lceil f_2/2\rceil\). Concavity and the integer end ordinates give
\[
 \lceil f_0/2\rceil+f_3+\beta_1+\beta_2-c_2\le3M+1.
\]
Indeed \((f_0+2f_3)/3\le f_2\) gives the first two terms at most \(3M\).
The indicator correction is at most 1 unless \(f_1,f_2\) are integers
and \(f_2=2M-1\). In that case, writing \(f_0=2a+p\), \(p=0,1\),
concavity gives \(f_1\ge a+M\) and \(f_3\le3M-a-2\), improving
the first bound to \(3M-1\).
Adding both chains and the extra point at the O end gives \(N\le3i+6\).

For \(W=4\), pattern EE gives
\[
 \lceil f_0/2\rceil+\lceil f_4/2\rceil+
 \beta_1+\beta_2+\beta_3-c_2\le2\lceil f_2/2\rceil+3.
\]
The end terms are at most \(f_2+1\). The indicator sum is at most 2,
unless all three are 1 and \(f_2\) is odd; then \(f_2+1=2\lceil f_2/2\rceil\).
Adding the chains gives \(N\le2i+8\le3i+6\).
For OO, concavity gives \(f_0+f_4\le f_1+f_3\), and
\(\beta_j-c_j\le2\lceil f_j/2\rceil-f_j\) on columns 1,3.
With \(\beta_2\le1\), adding the chains and the two O end points gives
the same \(N\le2i+8\).

For the wider cases, let \(n\) be the number of interior coarse columns.
Their lengths sum to at most \(2i+2n\). End-chord concavity and at most
two fine boundary points per interior fine column give
\[
\begin{array}{c|c|c}
\text{pattern}&W&\text{upper bound for }N\\ \hline
\mathrm{EE}&2n+2&2i/n+4n+6\\
\mathrm{OO}&2n&4i/n+4n+4\\
\mathrm{EO}&2n+1&(4n+2)i/[n(n+1)]+4n+6-2/(n+1).
\end{array}
\]
For EE/OO the coarse columns are symmetric, and their average length
dominates the average end length. An E end has at most \(\ell/2+1\)
nonoriginal points, an O end at most \(\ell+1\).
For EO, the coarse columns \(2,4,\ldots,2n\) instead give
\[
 \sum\ell_j\ge\frac{n^2}{2n+1}\ell_0+
                  \frac{n(n+1)}{2n+1}\ell_W.
\]
Multiplying by \((2n+1)/(n(n+1))\) bounds \(\ell_0/2+\ell_W\),
since \(n/(n+1)\ge1/2\).

For \(2\le n\le8\), all three bounds are at most \(3i+6\) when \(i\ge13\).
For \(n\ge9\), the minimizing section bound, in original units, gives
\[
 i\ge\sum_{j=1}^n(\min\{j,n+1-j\}-1)
   =\lfloor(n-1)^2/4\rfloor\ge(n^2-2n)/4.
\]
Substitution reduces the three comparisons to
\(3n^2-24n+4\ge0\), \(3n^2-26n+16\ge0\), and
\(3n^3-23n^2-16n+12\ge0\). Each is positive at 9 and increasing thereafter.
Widths below 3 have no coarse interior column or are the thin line case.
\end{proof}

\subsection{A joint bound at larger widths}
\begin{lemma}\label{lem:joint}
If \(i\ge17\) and the minimum lattice width of \(P\) exceeds 2, then
\(L_P(2)\le6i+2b+6\), equivalently \(A+B\le12i+4b+10\).
\end{lemma}
\begin{proof}
Use the minimizing fine projection and E/O convention just introduced.
Its patterns and coarse columns are
\[
\begin{array}{c|c|c|c}
 &W&\text{interior coarse columns}&\text{coarse ends}\\ \hline
\mathrm{EE}&2n+2&2,4,\ldots,2n&0,W\\
\mathrm{EO}&2n+1&2,4,\ldots,2n&0\\
\mathrm{OO}&2n&1,3,\ldots,2n-1&\text{none}.
\end{array}
\]
Here \(n\ge2\), and in the wide OO case \(n\ge3\).
Let \(p_j\) count all closed fine points on column \(j\), and let
\(i_j,b_j\) be its original interior and boundary counts on coarse columns.
Then \(\sum p_j=L_P(2)\), \(\sum i_j=i\), \(\sum b_j=b\).
All variables are nonnegative. Convexity and interval counting give
\begin{align}
 \ell_{j-1}-2\ell_j+\ell_{j+1}&\le0,\qquad \ell_j\ge\min\{W/2,j,W-j\},\label{eq:sectionLP}\\
 p_j&\le\ell_j+1 \qquad\text{on all columns},\nonumber\\
 \ell_j&\le2i_j+2,\quad 2i_j\le\ell_j+2 \qquad\text{on interior coarse columns},\nonumber\\
 p_j&\le2i_j+b_j+1,\quad b_j\le2
 \qquad\text{on interior coarse columns},\label{eq:coarseLP}\\
 \ell_j&\le2b_j,\quad p_j\le2b_j+1
 \qquad\text{on coarse ends}.\label{eq:endLP}
\end{align}
For the closed coarse bound, between two odd fine points is an even point
strictly inside the section, so there are at most \(i_j+1\) odd points.
At an end, the interval has integer endpoints; its even points give the
two end bounds even when it is a single odd vertex.

For \(2\le n\le8\), relax these conditions to nonnegative real variables and
add \(\sum i_j\ge17\). Write the displayed constraints as \(Mx\le r\), with
objective \(c x=\sum p_j-6\sum i_j-2\sum b_j\).
The companion file \texttt{dual\_weights.json} contains rational vectors
\(\lambda\ge0\) with
\[
 M^{\mathsf T}\lambda\ge c^{\mathsf T},\qquad
 r^{\mathsf T}\lambda=\gamma\le6.
\]
Their exact values are checked by the supplied rational verifier. The bounds are
\begin{center}
\begin{tabular}{c r r r}
\toprule
\(n\)&EE&EO&OO\\ \midrule
2&\(-4\)&\(-5\)&\(-6\)\\
3&\(-17/3\)&\(-20/3\)&\(-23/3\)\\
4&\(-9/2\)&\(-11/2\)&\(-13/2\)\\
5&\(-11/5\)&\(-16/5\)&\(-21/5\)\\
6&\(2/3\)&\(-1/3\)&\(-4/3\)\\
7&\(7/3\)&\(8/3\)&\(5/3\)\\
8&\(9/2\)&\(94/21\)&\(7/2\)\\ \bottomrule
\end{tabular}
\end{center}
For every feasible \(x\ge0\),
\(cx\le\lambda^{\mathsf T}Mx\le\lambda^{\mathsf T}r\).
Thus these certificates prove the bound for all \(i\ge17\) at the indicated
indices. No floating optimization result is required. The extra OO
\(n=2\) certificate is not needed for the wide case.

For \(n\ge9\), let \(I_j\) count strict fine points on a coarse interior
column; interval parity gives \(I_j\le2i_j+1\).
Take the hull of the integer endpoints on columns \(j-1,j+1\).
For a two-dimensional integer trapezoid of width 2, Pick gives
\[
 p_{j-1}+p_{j+1}=p_j(T)+I_j(T)+2\le p_j+I_j+2.
\]
This is the row identity of \cite[Lemma 6.5]{HHK}.
If the hull is a segment, its two end sections are single points, so the
inequality holds directly. The minimum-section bound ensures all adjacent
sections contain fine points.

Summing these inequalities and subtracting the twice-counted intervening
noncoarse columns gives
\[
 L_P(2)-6i-2b\le
 \begin{cases}
 5n+2-S_{\mathrm{EE}},&\mathrm{EE},\\
 5n+1-S_{\mathrm{EO}},&\mathrm{EO},\\
 5n-S_{\mathrm{OO}},&\mathrm{OO}.
 \end{cases}
\]
Here the first two \(S\)'s sum \(p_j\) over \(j=3,5,\ldots,2n-1\);
the third sums over \(j=2,4,\ldots,2n-2\).
Indeed each coarse interior contributes
\(2p_j+I_j+2\le6i_j+2b_j+5\), and each coarse end contributes one
additional constant.
Since \(p_j\ge\min\{j,W-j\}\), pairing columns gives
\[
 S_{\mathrm{EE}}\ge\frac{n^2+2n-3}{2},\quad
 S_{\mathrm{EO}}\ge\frac{n^2+n-2}{2},\quad
 S_{\mathrm{OO}}\ge\frac{n^2-1}{2}.
\]
The exact EE sum is \(n^2/2+n-1\) for even \(n\), the stated value for odd
\(n\); the EO sum is exact and the OO sum is \(\lfloor n^2/2\rfloor\).
The resulting excess bounds are
\[
 \frac{-n^2+8n+7}{2},\quad
 \frac{-n^2+9n+4}{2},\quad
 \frac{-n^2+10n+1}{2}.
\]
At \(n=9\) these are \(-1,2,5\), and their successive differences are
\((7-2n)/2,(8-2n)/2,(9-2n)/2\), all negative thereafter.
This closes the tail and the proof.
\end{proof}

\section{Narrow widths and the exceptional area correction}\label{sec:narrow}
Width below 1 has no original interior line. At width 1 the admissible case
is the line family. At width \(3/2\), write the doubled projection as
\([-2,1]\). There is one coarse interior row 0. The end and middle bounds
\[
 p_{-2}\le2b_{-2}+1,\quad p_0\le2i+b_0+1,\quad I_0\le2i+1,
 \quad p_{-1}+p_1\le p_0+I_0+2
\]
give \(L_P(2)\le6i+2b+6\). Width 2 can exceed this by one.

\subsection{The width-two equality family}
\begin{lemma}\label{lem:top}
At minimum width 2, the EE pattern satisfies \(L_P(2)\le6i+2b+7\).
Its equality signatures are exactly \eqref{eq:topset}, and all are realized
for every \(i\ge2\). The OO pattern satisfies
\[
 L_P(2)-6i-2b\le12-i-b.
\]
For \(i\ge17\), all excess-seven signatures therefore belong to
\eqref{eq:topset}.
\end{lemma}
\begin{proof}
In the EE pattern index the projection of \(Q\) by \(0,\ldots,4\).
All original interior points are on column 2. Interval parity and the row
hull inequality give
\[
 L_P(2)\le(2b_0+1)+(2b_4+1)+2(2i+b_2+1)+(2i+1)+2
       =6i+2b+7.
\]
Equality forces both ordinates at each end to be odd integers.
Translate vertically so the middle original interior points have fine
ordinates \(2,4,\ldots,2i\). Equality in the strict fine count then forces
\[
 0\le v_2<1,\qquad 2i+1<u_2\le2i+2.
\]
Put \(c_L=1\) if \(v_2=0\), otherwise 0, and \(c_U=1\) if
\(u_2=2i+2\), otherwise 0. Then \(b_2=c_L+c_U\).
Writing \(L_j=\lceil v_j\rceil,U_j=\lfloor u_j\rfloor\) for \(j=1,3\),
equality in the row hull bound is equivalent to
\begin{equation}\label{eq:topends}
 L_1+L_3=1-c_L,\qquad U_1+U_3=4i+3+c_U.
\end{equation}
The hull has half-integer middle endpoints, and equality forces its strict
and closed middle counts to equal those of \(Q\); these are the only
possible endpoint sums. Conversely these conditions give
\(p_1+p_3=4i+b_2+4\), hence equality.

Let \(\beta_U,\beta_L\) count integer upper/lower ordinates on columns 1,3,
and put \(\beta=\beta_U+\beta_L\).
If \(c_U=1\), concavity and \eqref{eq:topends} force both upper ordinates
to be integer. If \(c_U=0\) and neither were integer, the upper chain
would have no vertex on columns 1,2,3, so it would join its two odd end
ordinates by one edge, making \(u_2\) integral, a contradiction.
Thus \(\beta_U\ge1+c_U\), and likewise \(\beta_L\ge1+c_L\).
Consequently
\[
 d=\beta-b_2-2\in\{0,1,2\},\qquad
 B=2b+2+\beta-b_2=2b+4+d.
\]
Pick with the equality value of \(L_P(2)\) gives
\(A=12i+2b+8-d\).

For the boundary range, let \(g_U=1\) if \(\beta_U=2\), otherwise 0,
and define \(g_L\) similarly. Concavity and odd end rounding give
\[
 u_0+u_4\le4i+2+2c_U-2g_U,\qquad
 v_0+v_4\ge2-2c_L+2g_L.
\]
To see the improvement when \(g_U=1\), concavity gives
\(u_0\le(3U_1-U_3)/2\), \(u_4\le(3U_3-U_1)/2\).
Rounding down to odd integers loses 2 in total if \(U_1+U_3=4i+4\),
and 3 if the sum is \(4i+3\). The lower-chain rounding is the reverse.
It follows that \(b\le2i+2b_2-g_U-g_L\).
Since \(g_U+g_L\ge\beta-2\),
\[
 b+2d\le2i+2\beta-g_U-g_L-4\le2i+\beta-2\le2i+2.
\]

For attainment, choose parameters as follows:
\[
\begin{array}{c|cc|cc}
d&c_L&c_U&\text{left capacity}&\text{right capacity}\\ \hline
0&1&1&i&i\\
1&1&0&i-1&i\\
2&0&0&i-1&i-1 .
\end{array}
\]
Set \(s=2-d\), \(b_0=\min\{\text{left capacity},b-s\}\),
\(b_4=b-s-b_0\). The range in \eqref{eq:topset} ensures the capacities.
Take the lower chain
\[
 (0,1),(1,0),(3,1-c_L),(4,v_4)
\]
and upper chain
\[
 (0,1+2b_0),(1,U_1),(3,2i+2),(4,v_4+2b_4),
\]
where \(v_4=1\) for \(d\le1\), \(v_4=3\) for \(d=2\), and
\(U_1=2i+2\) for \(d=0\), \(U_1=2i+1\) otherwise.
Their slopes have the required monotonicity by the capacity bounds.
Repeated endpoints are omitted. The original counts are \(i,b\), the
conditions \eqref{eq:topends} hold, and \(\beta=4\), giving the desired
\(A,B\). The inner rectangle for \(d=0,1\) has height at least \(2i+1\).
For \(d=2\), the inner parallelogram has width
\(|2a+c|+(2i+1)|c|\) in direction \((a,c)\).
Any direction with \(c\ne0\) therefore has fine width at least 5, whereas
the horizontal direction has width 4. Thus the minimum original width is 2.

For OO, the original columns are 1,3. Row pairing gives
\[
 L_P(2)\le2p_1+I_1+2p_3+I_3+4-p_2
        \le6i+2b+10-p_2.
\]
The coarse points in a section span length at least
\(2(i_j+b_j-1)\), including the empty case as a vacuous lower bound.
Concavity gives \(\ell_2\ge i+b-2\), whence \(p_2\ge i+b-2\).
This proves the OO bound. For smaller \(i\), its only pairs with possible
excess above 6 are \((2,2),(2,3),(3,2)\), included in the complete finite
source comparison of Section~\ref{sec:computation}.
\end{proof}

\subsection{The correction at \texorpdfstring{\(b=2,B=4\)}{b=2, B=4}}
\begin{lemma}\label{lem:neararea}
If \(i\ge17\) and a minimizing original projection has at least two interior
integral rows, then \(A<12i+2b+6\).
\end{lemma}
\begin{proof}
Let these rows be \(1,\ldots,n\), \(n\ge2\), and put
\(\ell_y=\operatorname{length}(P\cap\{y=\text{constant}\})\).
The half-row bands centered at \(j=1,\ldots,n\) have normalized area at most
\(8\ell_j\le8(i_j+1)\). The two exterior half-bands have total area at most
\(2(b_0+b_{n+1}+\ell_{1/2}+\ell_{n+1/2})\).
Concavity gives
\(\ell_{1/2}+\ell_{n+1/2}\le\ell_k+\ell_{n-k+1}\).
Among the \(\lfloor n/2\rfloor\) disjoint pairs, one has interior count
at most \(i/\lfloor n/2\rfloor\). Hence
\begin{equation}\label{eq:neararea}
 A\le8i+8n+2b+\frac{2i}{\lfloor n/2\rfloor}+4.
\end{equation}
For \(n=2\) this is strictly below \(12i+2b+6\) when \(i\ge8\).
For \(n\ge3\), it suffices that
\[
 i>f(n)=\frac{4n^2-5n+1}{2n-4}.
\]
At \(n=3,\ldots,7\), the values are
\(11,45/4,38/3,115/8,81/5\), all below 17.
At \(n=8,9,10\) they are \(217/12,20,351/16\).
The only remaining \(i\ge17\) are therefore among 17,\ldots,21.
The complete maximal-polygon classification gives the global minimum-width
maxima \(7,15/2,8,15/2,8\) at those five \(i\), respectively
\cite[proof of Theorem 1.4, pp.\ 9--10]{Inequalities}.
They exclude the required \(n=8,9,10\).
Their catalogue scope is explained in Section~\ref{sec:computation}.

For \(n\ge11\), the earlier section sum gives \(i\ge(n^2-2n)/4>f(n)\):
the latter strict comparison is equivalent to
\(n^3-12n^2+14n-2>0\), which is 31 at 11 and has increment
\(3n^2-21n+3>0\) thereafter.
\end{proof}

\begin{lemma}\label{lem:B4}
If \(i\ge17,b=2,B=4\), then \(A\le12i+8\).
\end{lemma}
\begin{proof}
With at least two minimizing interior original rows,
Lemma~\ref{lem:neararea} and even \(A\) give the result.
Width 1 is the line case and gives \(A=4i+4\).
The remaining one-row patterns are EE of fine height 4 and EO of fine height 3.
Integral polygons cannot have \(b=2\).

For EE put the fine end rows at \(-2,2\), and let \(c=0,1,2\) be the
original boundary count on row 0. If \(H_-,H_+\) are the fine end lengths,
the budget of four fine boundary points gives \(H_-+H_+\le2-c\).
In original units, trapezoid integration and concavity give
\[
 A\le12\ell_0+2(\ell_{-1}+\ell_1).
\]
The original middle endpoints have fractional parts in
\(\{0,1/6,1/4,1/3,1/2,2/3,3/4,5/6\}\), by interpolating between fine
vertex rows \(-2,-1,1,2\). Thus \(\ell_0\le i+2/3+c/6\).
Equality in the area concavity step makes both chains affine across the
fine rows \(-1,0,1\).

If \(c=0,H_-+H_+\le1\), the bound is \(12i+9\), and evenness gives
\(12i+8\). If the sum is 2, all four fine points lie on the ends;
there is no interior-row vertex and \(A=8\).
If \(c=2\), both ends are points and the two other fine boundary points
are the original middle endpoints. No odd-row vertex remains, so
\(A=4(2i+2)=8i+8\).
If \(c=1,H_-+H_+=1\), again no odd-row vertex remains.
The opposite middle endpoint is a nonintegral fine half-integer;
the middle fine length is at most \(2i+3/2\), giving \(A\le8i+8\).

For \(c=1,H_-+H_+=0\), an additional middle fine boundary point must be
odd, giving \(A\le12i+6\). Otherwise exactly one fine boundary point
is on an odd row, and the general bound is \(12i+10\).
Equality would make the chain through the original middle point affine:
if it is integer on one odd row, it is integer on the other, requiring
two fine boundary points there. If neither is integer, it has no vertex
there or at the middle under this affineness condition. Its end points
have opposite fine horizontal parity, so their middle average is
nonintegral, contradicting the original middle point. Equality is impossible,
and evenness again gives \(A\le12i+8\).

For EO put the fine ends at \(-2,1\), let \(c\) be the middle original
boundary count and \(H\) the lower fine end length. The lower original
count is \(2-c\), and \(H\le2-c\).
If \(c=0\), then \(H=2\), all fine boundary is at the ends, and \(A=6\).
If \(c=2\), \(H=0\), the lower end is odd, the upper is a point, and
no vertex remains at row \(-1\); thus \(A=3(2i+2)\).
For \(c=1,H=1\), the same absence of an odd-row vertex gives
\(A=2+3H_0\le6i+7\), where \(H_0\le2i+5/3\) is the fine middle length.
For \(c=1,H=0\), a two-point upper end gives \(A=3H_0+1\le6i+6\).
If instead the upper end is a point and both middle endpoints are fine,
the nonoriginal one is odd, giving \(A=3H_0\le6i+3\).

The last EO case has four fine points: an original lower point \(E\),
an original middle point \(V\), an upper point \(T\) and a point \(W\)
on row \(-1\). An integer shear and translation put \(E=(0,-2),V=(0,0)\).
Reflect horizontally so \(T=(a,1)\), \(a>0\).
Convexity gives \(W=(c',-1)\), \(-a\le c'\le0\).
The middle fine interval is \([0,2a/3]\). The edge \(ET\) is primitive,
so \(3\nmid a\) and \(i=\lceil a/3\rceil-1\), whence
\(a=3i+1\) or \(3i+2\). Shoelace gives
\(A=2(a-c')\le4a\le12i+8\), finishing every case.
\end{proof}

For \(i\ge17\), the marginal bounds, Lemma~\ref{lem:lower},
Lemma~\ref{lem:boundarycap}, Lemma~\ref{lem:joint} and the narrow-width
analysis put every signature in the body, top or line.
The generic maximum-area correction is \eqref{eq:maxlower}.
At \(b=2\), Lemma~\ref{lem:B4} excludes \(B=4\) at the two even areas
\(12i+10,12i+12\), giving precisely the two corrections in \eqref{eq:bodyB}.
This proves necessity in Theorem~\ref{thm:signatures} for \(i\ge17\).

\section{Uniform witnesses}\label{sec:witnesses}
Every polygon in this section is \(P=Q/2\), where \(Q\) is the hull of
the listed integer points. Collinear or repeated outline points are removed.
For successive collinear primitive-direction subedges of lengths \(p,q\),
their boundary gcds add to \(p+q\); hence merging them does not alter the
boundary computations. A listed original boundary point need not remain
an extreme vertex. These conventions include all parameter endpoints.

\subsection{The line family}
For \(4i+4\le A\le8i+8\), set \(n=A-4i-4\),
\(r=\lfloor n/2\rfloor,s=\lceil n/2\rceil\), and take
\[
 Q=\conv\{(0,-1),(s,-1),(2i+2,0),(r,1),(0,1)\}.
\]
Since \(0\le r,s\le2i+2\), this is convex and two-dimensional.
The original middle section is \([0,i+1]\); its endpoints are the two
boundary points, with \(i\) interior points. Integration gives
\(A=r+s+4i+4\), and \(i(Q)=2i+1\) gives \(B=A-4i\) by Pick.
This realizes the entire line family. Its members are distinct by area.

\subsection{Body coordinates and the two outer layers}
For a body triple put
\[
 m=3i-b+6,\quad e=A\bmod2,\quad
 j=\frac{A-6i-4b+6-e}{2},\quad k=\frac{B-2b-e}{2}.
\]
Its conditions become
\begin{equation}\label{eq:bodycoords}
 j,k\ge0,\quad k\le j,\quad 2k+e\le m,\quad j+k+e\le m+2.
\end{equation}
For \(m\ge2\), the area bound gives \(j\le m+1\) at \(e=0\) and
\(j\le m\) at \(e=1\). It excludes
\((e,j,k)=(0,m+1,0)\) except at \(b=2i+4\), and also excludes
\((0,m,0)\) at \(b=2\). At \(m=0\) only \(j=k=e=0\) occurs;
at \(m=1\), \(k=0\) and \(j,e\in\{0,1\}\).
We first realize the layers \(j+k+e=m+2,m+1\).

If \(j+k+e=m+2\), then \(k\ge1\). Put \(f=2k+e-2\) and take
\begin{equation}\label{eq:outer2}
 Q=\conv\{(0,1),(f,1),(4i+4-f,-1),(2b-3,-2),(1,-2),(0,-1)\}.
\end{equation}
Here \(0\le f\le m-2\).
Right-chain convexity requires \(2b-3\le6i+6-2f\), which follows
strictly from this bound. The left chain is convex as well.
The original bottom has \(b-2\) boundary points; the middle fine
section is \([0,2i+2]\), giving \(i\) interior and two original boundary
points. On fine rows \(-2,-1,0,1\), the closed counts are
\[
 2b-3,\quad4i+5-f,\quad2i+3,\quad f+1.
\]
Thus \(L_P(2)=6i+2b+6\), \(B=2b+f+2=2b+2k+e\), and
\[
 A=2L_P(2)-B-2=12i+2b+10-2k-e
   =6i+4b-6+2j+e.
\]

On \(j+k+e=m+1\), first suppose \(b\ge3\), \(2k+e\ge1\).
With \(f=2k+e-1\), take
\begin{equation}\label{eq:outer1}
 Q=\conv\{(0,1),(f,1),(4i+4-f,-1),(2b-5,-2),(0,-2)\}.
\end{equation}
Its mixed-parity bottom has \(b-2\) original boundary points, the middle
has \(i\) interior and two boundary points, and right convexity follows
from \(f\le m-1\). The four fine row counts are
\(2b-4,4i+5-f,2i+3,f+1\). Hence \(L_P(2)=6i+2b+5\),
\(B=2b+f+1\), and Pick gives the target \(A\).

For \(b=2,k\ge1,2k+e\le m-1\), put \(f=2k+e-2\) and take
\[
 Q=\conv\{(0,1),(f,1),(4i+3-f,-1),(1,-2),(0,-2)\}.
\]
The bottom contains one original boundary point and the middle fine
section is \([0,(4i+3)/2]\), with \(i\) original interior points and
one original boundary point. The condition \(f\le3i+1\) gives convexity.
The fine counts \(2,4i+4-f,2i+2,f+1\) give
\(L_P(2)=6i+9,\ B=f+6=4+2k+e\), and the required area.

At \(b=2,2k+e=m\), necessarily \(j=k+1\); use
\begin{equation}\label{eq:b2cap}
 Q=\conv\{(0,1),(3i,1),(3,-2),(1,-2),(0,-1)\}.
\end{equation}
Its middle fine section is \([0,2i+1]\), and the original bottom point
is \((2,-2)\) in fine coordinates. Thus the counts are \(i,b=2\).
The edge gcds are \(3i,3,2,1,2\), so \(B=3i+8\);
shoelace gives \(A=9i+8\), exactly the desired signature.

For \(e=k=0\), the first outer layer is allowed only at \(b=2i+4\).
The triangle
\[
 Q=\conv\{(-1,-2),(4i+3,-2),(1,2)\}
\]
has \(i\) original interior points, \(b=2i+4\), \(B=4i+8\),
\(A=16i+16\). It supplies that exception.
The remaining first-layer case \(b=2,e=1,k=0\) is covered by the
high-difference construction below.

\subsection{The bulk}
It remains to cover \(j+k+e\le m\).
For \(b\ge3\), if \(e=0\) or \(e=1,j\ge k+1\), take
\[
 Q=\conv\{(0,1),(2k,1),(2i+2,0),
          (i+b-2+j-k,-1),(2b-6+e,-2),(0,-2)\}.
\]
The right-middle cross products are \(j+k-m\) and \(e-2(j-k)\),
both nonpositive; the other corners have the required sign.
The original middle section is \([0,i+1]\), and the bottom contains
\(b-2\) original boundary points.
Edge gcds give \(B=2b+2k+e\); shoelace gives
\(A=6i+4b-6+2j+e\).
For \(e=1,j=k\), instead take
\[
 Q=\conv\{(0,1),(2k+1,1),(2i+2,0),(2b-6,-2),(0,-2)\}.
\]
Its right convexity condition is \(2k+1\le m\), and the original counts
are unchanged. Here \(B=2b+2k+1\), \(A=6i+4b-5+2k\).
These formulas include the allowed \(m=0,1\) cases.

For \(b=2,j-k\ge2,2k+e\ge1\), use
\begin{equation}\label{eq:highdiff}
 Q=\conv\{(0,1),(2k+e-1,1),(2i+2,0),
          (i+j-k,-1),(1,-2),(0,0)\}.
\end{equation}
It is convex whenever \(j+k+e\le m+1\) and \(3-2(j-k)\le0\),
covering both its bulk uses and the remaining first-layer case.
The fine middle endpoints are \(0,2i+2\); the odd bottom vertex
has no original point. Thus the original counts are \(i,2\).
Gcds and shoelace give \(B=4+2k+e\), \(A=6i+2+2j+e\).

For \(b=2,e=k=0\), the exclusions give \(0\le j\le3i+3\).
Put \(h=3i+1,H=0\) at \(j=0\), and \(h=3i+2,H=j-1\) at \(j\ge1\).
Take
\begin{equation}\label{eq:B4witness}
 Q=\conv\{(0,0),(1,H),(2,0),(3,-h)\}.
\end{equation}
The condition \(H\le h\) gives convexity, including the collinear endpoint.
The upper original triangle lies between \(x=0,1\), with apex at \(x=1/2\),
so it has no original interior point.
On original row \(y=-r\), the lower triangle has endpoints
\(3r/h,1+r/h\). Its only possible original interior coordinate is \(x=1\),
which is inside exactly for \(0<r<h/3\). Hence its count is
\(\lceil h/3\rceil-1=i\).
The two original boundary points are the base endpoints. Every nonbase
outline subedge is primitive, since \(\gcd(3,h)=1\), giving \(B=4\);
when \(H=0\), the base has fine length 2 and gives the same count.
The area is \(A=2(h+H)=6i+2+2j\).
At \(H=h\), the fine point \((2,0)\), corresponding to an original lattice
point, lies on a merged edge rather than remaining a vertex. Its boundary
status is unchanged.
In particular, this body family includes \(A=8i,B=4\), satisfying
\eqref{eq:collapse}.

Finally let \(b=2,\delta=j-k\in\{0,1\}\), \(t=2k+e\le3i+4\).
If \(t\le3i+1+\delta\), put \(D=6i+2+2\delta-2t\) and take
\[
 Q=\conv\{(0,0),(0,1),(t,1),(D,-2)\}.
\]
Here \(D\) is even and nonnegative, giving convexity.
Its original middle interval is
\([0,i+(1+\delta)/3]\), so it has \(i\) interior points and one
boundary point there, and the bottom supplies the other boundary point.
The right edge is primitive because \(D-t\) is 2 or 1 modulo 3.
Thus \(B=t+4\), and \(A=3t+D=6i+2+t+2\delta\).

For \(\delta=0,t=3i+4\), use
\[
 Q=\conv\{(0,1),(3i-2,1),(7,-2),(3,-2)\}.
\]
Its fine middle interval is \([1,2i+1]\); the bottom odd-ended interval
has two original points. The edge gcds give \(B=3i+8\), and
\(A=9i+6\).
For \(t=3i+3\), replace the upper-left corner by \((1,1),(1,0)\).
For \(t=3i+2\), also replace the upper-right corner by
\((3i-3,1),(2i+1,0)\).
Each removed primitive triangle has normalized area 1; replacing its two
primitive old edges by one reduces \(B\) by 1 as well.
All removed and new fine points are nonoriginal, so \(i,b\) do not change.
The top segment stays nonempty at \(i=2\).

For \(\delta=1,t=3i+4\), use \eqref{eq:b2cap}.
For \(t=3i+3\), shave its upper-right primitive corner, replacing
\((3i,1)\) by \((3i-1,1),(2i+1,0)\).
Again \(A,B\) both decrease by 1 without changing \(i,b\).
Together these families exhaust \eqref{eq:bodycoords}.
Thus every body signature is realized for every \(i\ge2\).
The top witnesses were already constructed in Lemma~\ref{lem:top}.

\section{Computational inputs and reproduction}\label{sec:computation}
The complete finite input is the published classification
\cite[Classification 6.7 and the following enumeration]{Classification},
with the released dataset \cite{Dataset}. Its convention is \(k\)-rational,
meaning \(kP\) is integral, so \(k=2\) includes integral polygons.
Affine equivalence uses integer translations and unimodular matrices; these
preserve \(L_P(t)\) for every integer \(t\).

We explain the completeness and count parameters used here.
Every positive-\(i\) half-integral polygon is contained in a same-\(i\)
2-maximal polygon. Indeed the minimizing section estimate gives, for every positive \(i\),
\[
 i\ge\lfloor(n-1)^2/4\rfloor\ge(n^2-2n)/4
\]
when there are \(n\) interior integral columns. Since \(w\le n+1\), this
implies \(w\le2+\sqrt{4i+1}\). Some integral section has length at most
\(i+1\). A chord to a longest section, at distance at least \(1/2\) from
the projection end, gives \(\ell_{\max}\le2w(i+1)\).
Thus \(\area(P)\le2w^2(i+1)\), a bound depending only on \(i\).
The integer normalized area increases under proper half-integral containment,
so a same-\(i\) enlargement chain terminates.
Successively deleting a fine vertex outside a target subpolygon reaches
every subpolygon; along a same-\(i\) chain no interior count is lost.
This is the shaving principle in \cite[Proposition 2.9, Algorithm 1 and
Remark 2.14]{Classification}.

The maximal catalogue combines strip polygons with Fine interiors.
For an integral polytope \(K\), its Fine interior is
\[
 F(K)=\{x:\langle x,u\rangle\ge\min_{y\in K}\langle y,u\rangle+1
               \text{ for every }u\in\Z^{\dim K}\setminus\{0\}\}.
\]
Its integer points are exactly the interior integer points of \(K\).
For a half-integral polygon \(P_0\), the lattice pyramid with base
\(2P_0\) at height \(-1\) and apex \((0,0,1)\) has middle section \(P_0\).
Its interior integer points lie at height 0 and number \(i(P_0)\).
Positive interior count forces every lattice width to be at least 2, and the
height direction has width 2. Its Fine interior lies in the middle plane.
Consequently its Fine polygon has exactly \(i(P_0)\) total original lattice
points, including its own boundary.

The source results \cite[Propositions 6.4 and 6.5]{Classification} and
\cite[Theorems 2.5, 3.2 and 4.1; Remark 3.16]{Fine} put the
one-dimensional Fine case in the strip branch and reconstruct the
two-dimensional case from the complete Fine catalogue through 40 total points.
Hence the maximal catalogue through 40 interior points covers all arrangements,
not only collinear ones. The same-\(i\) containment argument and monotonicity
of minimum width justify the five width maxima used in
Lemma~\ref{lem:neararea}.

For \(i=2,\ldots,16\), version 1.0.0 of the dataset provides the full function
files \texttt{ehrhart\_iN.h5}; its separate collinear polygon files are not
used. Each normalized coefficient matrix is stored as an integer
\(3\times2\) array
\[
 \begin{pmatrix}A&A\\4b&2B\\8i+4b-A&8\end{pmatrix}.
\]
The decoder verifies all entries, the area group, parity and uniqueness,
then filters only \(b\ge2\). The resulting exact sets \(S_i\) have empty
symmetric difference with \(C_i\) at every one of these fifteen \(i\).
For reference, their cardinalities are
\[
\begin{array}{c|rrrrrrrr}
i&2&3&4&5&6&7&8&9\\ \hline
|S_i|&408&761&1267&1953&2846&3973&5361&7037\\
\end{array}
\]
\[
\begin{array}{c|rrrrrrr}
i&10&11&12&13&14&15&16\\ \hline
|S_i|&9028&11361&14063&17161&20682&24653&29101 .
\end{array}
\]
The certificate is equality of the full sets, not just these counts.
Combined with the published completeness assertion this proves finite
necessity. With the analytic necessity for \(i\ge17\) and the constructions
in Section~\ref{sec:witnesses}, it completes Theorem~\ref{thm:signatures}.

The companion Python source includes the exact section model and its 21
rational dual weights, a verifier using integer and rational arithmetic,
the source-file checksum manifest and coefficient-set comparator, and the
line/body/top constructors. Polygon counts use integer half-plane
intersection, shoelace area and edge gcds; small cases are also checked by
independent bounding-box enumeration. Data retrieval is pinned to version 1.0.0
and its official checksums. Exact source-set and inequality verification can
be reproduced without rerunning the polygon classification.
The companion archive \texttt{half-integral-ehrhart-verification.zip}
includes a root \texttt{README.md} with fixed-input retrieval,
full-comparison commands and pinned optional HDF5 dependencies.
From the archive root, run
\verb|python3 numerical_case/verify.py --output dual.json|
for the certificate-only check.

The published catalogue computations are cited external inputs and have not
been independently reconstructed here. Floating linear optimization was used
only to discover rational weights; the retained rational inequalities suffice
for their verification. Finite polygon tests check the implementation, whereas
the parameter arguments in the preceding sections establish universal
attainment.

\section{The exact count}\label{sec:count}
For a body triple recall \(m=3i-b+6\) and the coordinates
\eqref{eq:bodycoords}. As \(b\) runs from 2 to \(3i+6\), \(m\) runs from
\(N=3i+4\) to 0. Temporarily remove the generic maximum-area point
\((e,j,k)=(0,m+1,0)\). For \(m\ge2\), at \(e=0,k=0\) there are \(m+1\)
choices, and at \(e=0,k\ge1\) there are \(m+3-2k\), for
\(1\le k\le\lfloor m/2\rfloor\).
At \(e=1,k=0\) there are \(m+1\) choices, and at \(e=1,k\ge1\) there
are \(m+2-2k\), for \(1\le k\le\lfloor(m-1)/2\rfloor\).
Summing gives \(m(m+7)/2\) for both parities.
The \(m=1\) count is 4 and \(m=0\) count is 1.
Reinstating the exceptional point at \(b=2i+4\) adds one, while the
additional \(b=2,j=m,k=0,e=0\) exclusion removes one. Therefore
\[
 \#\text{body}
 =1+\sum_{m=1}^{N}\frac{m(m+7)}2
 =1+\frac{N(N+1)(N+11)}6
 =\frac92i^3+36i^2+\frac{155}{2}i+51.
\]
The top count is
\((2i+1)+(2i-1)+(2i-3)=6i-3\), and the line count is \(4i+5\).
These sets are disjoint: the body has \(A+B\le12i+4b+10\),
the top has \(A+B=12i+4b+12\), and their respective \(A-B\) values
are at least \(6i+2b-6>4i\) and \(12i+4-2d>4i\),
whereas the line has \(A-B=4i\).

\begin{proof}[Proof of Theorem~\ref{thm:count}]
Theorem~\ref{thm:signatures} identifies the attained signatures with \(C_i\).
The constituent formula \eqref{eq:coeff} identifies signatures with entire
functions at fixed \(i\). Adding the three disjoint counts gives
\[
 \Ehr(i)=\frac92i^3+36i^2+\frac{175}{2}i+53
        =\frac{9i^3+72i^2+175i+106}{2}.
\]
All parameter ranges include \(i=2\), integral polygons and period collapse.
\end{proof}

\paragraph{AI-use disclosure.}
OpenAI Codex assisted with proof development, exact verification code and
manuscript drafting. The cited computational classifications remain external
dependencies. The mathematical arguments and accompanying verification inputs
are supplied for independent checking.

\begin{thebibliography}{9}
\bibitem{Scott}
P.~R. Scott, \emph{On convex lattice polygons},
Bulletin of the Australian Mathematical Society \textbf{15} (1976), 395--399.
\href{https://doi.org/10.1017/S0004972700022826}{doi:10.1017/S0004972700022826}.
\bibitem{Collapse}
M. Bohnert, \emph{Quasi-period collapse in half-integral polygons},
arXiv:2405.13404v2, 2 July 2025.
\href{https://arxiv.org/abs/2405.13404v2}{arXiv:2405.13404v2}.
\bibitem{Inequalities}
M. Bohnert and J. Springer,
\emph{Generalizations of Scott's inequality and Pick's formula to rational polygons},
arXiv:2411.11187v1, 17 November 2024.
\href{https://arxiv.org/abs/2411.11187v1}{arXiv:2411.11187v1}.
\bibitem{Classification}
M. Bohnert and J. Springer,
\emph{Classifying Rational Polygons with Small Denominator and Few Interior Lattice Points},
Discrete \& Computational Geometry \textbf{76} (2026), 1238--1265.
\href{https://doi.org/10.1007/s00454-026-00822-0}{doi:10.1007/s00454-026-00822-0}.
\bibitem{AreaBounds}
M. Bohnert,
\emph{Area bounds for planar convex bodies containing a fixed number of interior integral points},
arXiv:2305.11485v2, 2 July 2025.
\href{https://arxiv.org/abs/2305.11485v2}{arXiv:2305.11485v2}.
\bibitem{Fine}
M. Bohnert, \emph{Lattice 3-polytopes of lattice width 2 and corresponding toric hypersurfaces},
arXiv:2412.17545v2, 3 July 2025.
\href{https://arxiv.org/abs/2412.17545v2}{arXiv:2412.17545v2}.
\bibitem{HHK}
G. Hamm, J. Hofscheier and A. Kasprzyk,
\emph{Classification and Ehrhart theory of denominator 2 polygons},
arXiv:2411.19183v1, 28 November 2024.
\href{https://arxiv.org/abs/2411.19183v1}{arXiv:2411.19183v1}.
\bibitem{Dataset}
J. Springer and M. Bohnert,
\emph{Half-integral polygons with few interior lattice points},
dataset, version 1.0.0, 14 October 2024.
\href{https://doi.org/10.5281/zenodo.13928298}{doi:10.5281/zenodo.13928298}.
\end{thebibliography}
\end{document}
